
\documentclass{article}
\usepackage{colortbl}
\usepackage{makecell}
\usepackage{multirow}
\usepackage{supertabular}

\begin{document}

\newcounter{utterance}

\centering \large Interaction Transcript for game `hot\_air\_balloon', experiment `air\_balloon\_survival\_ro\_complexity\_easy', episode 1 with Qwen3.6{-}35B{-}A3B{-}FP8{-}without{-}reasoning.
\vspace{24pt}

{ \footnotesize  \setcounter{utterance}{1}
\setlength{\tabcolsep}{0pt}
\begin{supertabular}{c@{$\;$}|p{.15\linewidth}@{}p{.15\linewidth}p{.15\linewidth}p{.15\linewidth}p{.15\linewidth}p{.15\linewidth}}
    \# & $\;$A & \multicolumn{4}{c}{Game Master} & $\;\:$B\\
    \hline

    \theutterance \stepcounter{utterance}  
    & & \multicolumn{4}{p{0.6\linewidth}}{
        \cellcolor[rgb]{0.9,0.9,0.9}{
            \makecell[{{p{\linewidth}}}]{
                \texttt{\tiny{[P1$\langle$GM]}}
                \texttt{Participați la un joc de negociere colaborativă.} \\
\\ 
\texttt{Împreună cu un alt participant, trebuie să conveniți asupra unui singur set de elemente care vor fi păstrate. Fiecare dintre voi are propria viziune asupra importanței fiecărui element pentru sine. Nu știți cum evaluează celălalt participant elementele. În plus, vi se oferă efortul pe care îl necesită fiecare element.} \\
\texttt{Puteți conveni asupra unui set numai dacă efortul total al elementelor selectate nu depășește o limită comună:} \\
\\ 
\texttt{LIMITĂ = 5174} \\
\\ 
\texttt{Iată valorile efortului individual pentru fiecare element:} \\
\\ 
\texttt{Efort element = \{"C08": 400, "C55": 990, "A43": 315, "A67": 43, "B42": 336, "A62": 192, "B19": 325, "C43": 812, "C73": 868, "C10": 593, "B36": 915, "A18": 944, "C61": 311, "C67": 252, "C42": 343, "C70": 104, "B40": 558, "C39": 627, "A21": 593, "C16": 827\}} \\
\\ 
\texttt{Iată viziunea dvs. personală asupra importanței fiecărui element:} \\
\\ 
\texttt{Valori importanță element = \{"C08": 138, "C55": 583, "A43": 868, "A67": 822, "B42": 783, "A62": 65, "B19": 262, "C43": 121, "C73": 508, "C10": 780, "B36": 461, "A18": 484, "C61": 668, "C67": 389, "C42": 808, "C70": 215, "B40": 97, "C39": 500, "A21": 30, "C16": 915\}} \\
\\ 
\texttt{Obiectiv:} \\
\\ 
\texttt{Obiectivul dvs. este de a negocia un set comun de elemente care să vă avantajeze cât mai mult posibil (adică să maximizeze importanța totală pentru DUMNEAVOASTRĂ), respectând în același timp LIMITA. Nu sunteți obligat să faceți o PROPUNERE în fiecare mesaj {-} puteți pur și simplu să negociați. Toate tacticile sunt permise!} \\
\\ 
\texttt{Protocol de interacțiune:} \\
\\ 
\texttt{Puteți utiliza doar următoarele formate structurate într{-}un mesaj:} \\
\\ 
\texttt{PROPUNERE: \{'A', 'B', 'C', …\}} \\
\texttt{Propuneți păstrarea exact a acestor elemente.} \\
\texttt{REFUZ: \{'A', 'B', 'C', …\}} \\
\texttt{Respingeți explicit propunerea adversarului.} \\
\texttt{ARGUMENT: \{'...'\}} \\
\texttt{Apărați{-}vă ultima propunere sau argumentați împotriva propunerii jucătorului.} \\
\texttt{ACORD: \{'A', 'B', 'C', …\}} \\
\texttt{Acceptați propunerea adversarului, ceea ce încheie jocul.} \\
\\ 
\\ 
\texttt{Reguli:} \\
\\ 
\texttt{Puteți da ACORD doar pentru o propunere pe care cealaltă parte a înregistrat{-}o prin PROPUNERE.} \\
\texttt{Puteți da REFUZ doar pentru o propunere pe care cealaltă parte a înregistrat{-}o prin PROPUNERE.} \\
\texttt{Efortul total al oricărui set din PROPUNERE sau ACORD trebuie să fie ≤ LIMITĂ.} \\
\texttt{NU vă dezvăluiți scorurile de importanță ascunse.} \\
\texttt{O etichetă într{-}un format structurat trebuie să fie urmată de două puncte și spațiu alb. Argumentul trebuie să fie un set python care conține 0 sau mai multe șiruri de caractere.} \\
\texttt{Deci, trebuie să fie de forma "ETICHETĂ: \{...\}"} \\
\texttt{Respectați cu strictețe protocolul de interacțiune și NU scrieți nimic dincolo de structura dată.} \\
\texttt{Jocul se încheie atunci când o parte dă un ACORD la o PROPUNERE făcută de celălalt jucător.} \\
\texttt{Conținutul din răspunsul dvs. care poate fi transmis celuilalt jucător trebuie să fie non{-}gol.} \\
\texttt{Sunt active numai propunerile care au fost înregistrate prin structura de format PROPUNERE și care nu au fost refuzate prin REFUZ.} \\
\texttt{Trebuie să includeți formatul ARGUMENT cel puțin o dată undeva în toate mesajele dumneavoastră.} \\
            }
        }
    }
    & \\ \\

    \theutterance \stepcounter{utterance}  
    & \multicolumn{4}{p{0.6\linewidth}}{
        \cellcolor[rgb]{0.9,0.9,0.9}{
            \makecell[{{p{\linewidth}}}]{
                \texttt{\tiny{[P1$\rangle$GM]}}
                \texttt{ARGUMENT: \{'Încep negocierile. Scopul nostru este să maximizăm valoarea totală pentru amândoi, respectând limita de efort de 5174. Deoarece nu cunoaștem valorile celuilalt, propun o abordare bazată pe eficiență (raport valoare/efort). Vă rog să îmi transmiteți o propunere inițială sau să discutați despre criteriile de selecție.'\}} \\
            }
        }
    }
    & & \\ \\

    \theutterance \stepcounter{utterance}  
    & & \multicolumn{4}{p{0.6\linewidth}}{
        \cellcolor[rgb]{0.9,0.9,0.9}{
            \makecell[{{p{\linewidth}}}]{
                \texttt{\tiny{[GM$\rangle$P2]}}
                \texttt{Participați la un joc de negociere colaborativă.} \\
\\ 
\texttt{Împreună cu un alt participant, trebuie să conveniți asupra unui singur set de elemente care vor fi păstrate. Fiecare dintre voi are propria viziune asupra importanței fiecărui element pentru sine. Nu știți cum evaluează celălalt participant elementele. În plus, vi se oferă efortul pe care îl necesită fiecare element.} \\
\texttt{Puteți conveni asupra unui set numai dacă efortul total al elementelor selectate nu depășește o limită comună:} \\
\\ 
\texttt{LIMIT = 5174} \\
\\ 
\texttt{Iată valorile efortului individual pentru fiecare element:} \\
\\ 
\texttt{Item effort = \{"C08": 400, "C55": 990, "A43": 315, "A67": 43, "B42": 336, "A62": 192, "B19": 325, "C43": 812, "C73": 868, "C10": 593, "B36": 915, "A18": 944, "C61": 311, "C67": 252, "C42": 343, "C70": 104, "B40": 558, "C39": 627, "A21": 593, "C16": 827\}} \\
\\ 
\texttt{Iată viziunea dvs. personală asupra importanței fiecărui element:} \\
\\ 
\texttt{Item importance values = \{"C08": 138, "C55": 583, "A43": 868, "A67": 822, "B42": 783, "A62": 65, "B19": 262, "C43": 121, "C73": 508, "C10": 780, "B36": 461, "A18": 484, "C61": 668, "C67": 389, "C42": 808, "C70": 215, "B40": 97, "C39": 500, "A21": 30, "C16": 915\}} \\
\\ 
\texttt{Obiectiv:} \\
\\ 
\texttt{Obiectivul dvs. este de a negocia un set comun de elemente care să vă avantajeze cât mai mult posibil (adică să maximizeze importanța totală pentru DUMNEAVOASTRĂ), respectând în același timp LIMIT. Nu sunteți obligat să faceți o PROPUNERE în fiecare mesaj {-} puteți pur și simplu să negociați. Toate tacticile sunt permise!} \\
\\ 
\texttt{Protocol de interacțiune:} \\
\\ 
\texttt{Puteți utiliza doar următoarele formate structurate într{-}un mesaj:} \\
\\ 
\texttt{PROPUNERE: \{'A', 'B', 'C', …\}} \\
\texttt{Propuneți păstrarea exact a acestor elemente.} \\
\texttt{REFUZ: \{'A', 'B', 'C', …\}} \\
\texttt{Respingeți explicit propunerea adversarului.} \\
\texttt{ARGUMENT: \{'...'\}} \\
\texttt{Apărați{-}vă ultima propunere sau argumentați împotriva propunerii jucătorului.} \\
\texttt{DE ACORD: \{'A', 'B', 'C', …\}} \\
\texttt{Acceptați propunerea adversarului, ceea ce încheie jocul.} \\
\\ 
\\ 
\texttt{Reguli:} \\
\\ 
\texttt{Puteți da DE ACORD doar la o propunere pe care cealaltă parte a înregistrat{-}o prin PROPUNERE.} \\
\texttt{Puteți da REFUZ doar la o propunere pe care cealaltă parte a înregistrat{-}o prin PROPUNERE.} \\
\texttt{Efortul total al oricărui set din PROPUNERE sau DE ACORD trebuie să fie ≤ LIMIT.} \\
\texttt{NU vă dezvăluiți scorurile de importanță ascunse.} \\
\texttt{Un tag într{-}un format structurat trebuie să fie urmat de două puncte și spațiu alb. Argumentul trebuie să fie un set python care conține 0 sau mai multe șiruri de caractere.} \\
\texttt{Deci, trebuie să fie de forma "TAG: \{...\}"} \\
\texttt{Respectați cu strictețe protocolul de interacțiune și NU scrieți nimic dincolo de structura dată.} \\
\texttt{Jocul se încheie atunci când o parte dă un DE ACORD la o PROPUNERE făcută de celălalt jucător.} \\
\texttt{Conținutul din răspunsul dvs. care poate fi transmis celuilalt jucător trebuie să fie non{-}gol.} \\
\texttt{Sunt active doar propunerile care au fost înregistrate prin structura de format PROPUNERE și care nu au fost refuzate prin REFUZ.} \\
\texttt{Trebuie să includeți formatul ARGUMENT cel puțin o dată undeva în toate mesajele dumneavoastră.} \\
\\ 
\\ 
\texttt{Veți primi acum primul mesaj al celuilalt jucător.} \\
\\ 
\texttt{ARGUMENT: \{'Încep negocierile. Scopul nostru este să maximizăm valoarea totală pentru amândoi, respectând limita de efort de 5174. Deoarece nu cunoaștem valorile celuilalt, propun o abordare bazată pe eficiență (raport valoare/efort). Vă rog să îmi transmiteți o propunere inițială sau să discutați despre criteriile de selecție.'\}} \\
            }
        }
    }
    & \\ \\

    \theutterance \stepcounter{utterance}  
    & & & \multicolumn{4}{p{0.6\linewidth}}{
        \cellcolor[rgb]{0.9,0.9,0.9}{
            \makecell[{{p{\linewidth}}}]{
                \texttt{\tiny{[GM$\langle$P2]}}
                \texttt{ARGUMENT: \{'Am înțeles abordarea bazată pe eficiență. Pentru a începe, am calculat rapoartele valoare/efort pentru elementele cu cea mai mare importanță pentru mine. De exemplu, C16 are un raport excelent (915/827 ≈ 1.11), iar A67 este foarte eficient (822/43 ≈ 19.1). Propun să începem cu o selecție care să includă elementele cu cele mai mari rapoarte, pentru a asigura un echilibru bun între efort și valoare pentru ambele părți, deși voi prioritiza elementele cu impact major asupra scorului meu.'\}} \\
            }
        }
    }
     \\ \\

    \theutterance \stepcounter{utterance}  
    & & \multicolumn{4}{p{0.6\linewidth}}{
        \cellcolor[rgb]{0.9,0.9,0.9}{
            \makecell[{{p{\linewidth}}}]{
                \texttt{\tiny{[P1$\langle$GM]}}
                \texttt{ARGUMENT: \{'Am înțeles abordarea bazată pe eficiență. Pentru a începe, am calculat rapoartele valoare/efort pentru elementele cu cea mai mare importanță pentru mine. De exemplu, C16 are un raport excelent (915/827 ≈ 1.11), iar A67 este foarte eficient (822/43 ≈ 19.1). Propun să începem cu o selecție care să includă elementele cu cele mai mari rapoarte, pentru a asigura un echilibru bun între efort și valoare pentru ambele părți, deși voi prioritiza elementele cu impact major asupra scorului meu.'\}} \\
            }
        }
    }
    & \\ \\

    \theutterance \stepcounter{utterance}  
    & \multicolumn{4}{p{0.6\linewidth}}{
        \cellcolor[rgb]{0.9,0.9,0.9}{
            \makecell[{{p{\linewidth}}}]{
                \texttt{\tiny{[P1$\rangle$GM]}}
                \texttt{ARGUMENT: \{'Sunt de acord că A67 este extrem de eficient (822/43) și ar trebui să fie prioritar. Totuși, C16 are un raport de doar \textasciitilde{}1.11, ceea ce este relativ scăzut comparativ cu alte opțiuni. De exemplu, C42 are o valoare de 808 și un efort de 343 (raport \textasciitilde{}2.35), iar A43 are o valoare de 868 și un efort de 315 (raport \textasciitilde{}2.75). Acestea sunt mult mai eficiente decât C16. Propun să excludem C16 din setul inițial pentru a elibera efort pentru elemente cu randament mai mare. Vă propun următorul set inițial bazat pe eficiență ridicată pentru mine:'\}} \\
\texttt{PROPUNERE: \{'A67', 'A43', 'C42', 'C61', 'C10', 'B42', 'C73', 'C39', 'A18', 'B36', 'C55', 'C43', 'C70', 'B19', 'C67', 'C08', 'B40', 'A21'\}} \\
            }
        }
    }
    & & \\ \\

    \theutterance \stepcounter{utterance}  
    & & \multicolumn{4}{p{0.6\linewidth}}{
        \cellcolor[rgb]{0.9,0.9,0.9}{
            \makecell[{{p{\linewidth}}}]{
                \texttt{\tiny{[GM$\rangle$P2]}}
                \texttt{ARGUMENT: \{'Sunt de acord că A67 este extrem de eficient (822/43) și ar trebui să fie prioritar. Totuși, C16 are un raport de doar \textasciitilde{}1.11, ceea ce este relativ scăzut comparativ cu alte opțiuni. De exemplu, C42 are o valoare de 808 și un efort de 343 (raport \textasciitilde{}2.35), iar A43 are o valoare de 868 și un efort de 315 (raport \textasciitilde{}2.75). Acestea sunt mult mai eficiente decât C16. Propun să excludem C16 din setul inițial pentru a elibera efort pentru elemente cu randament mai mare. Vă propun următorul set inițial bazat pe eficiență ridicată pentru mine:'\}} \\
\texttt{PROPUNERE: \{'A67', 'A43', 'C42', 'C61', 'C10', 'B42', 'C73', 'C39', 'A18', 'B36', 'C55', 'C43', 'C70', 'B19', 'C67', 'C08', 'B40', 'A21'\}} \\
            }
        }
    }
    & \\ \\

    \theutterance \stepcounter{utterance}  
    & & & \multicolumn{4}{p{0.6\linewidth}}{
        \cellcolor[rgb]{0.9,0.9,0.9}{
            \makecell[{{p{\linewidth}}}]{
                \texttt{\tiny{[GM$\langle$P2]}}
                \texttt{REFUZ: \{'A67', 'A43', 'C42', 'C61', 'C10', 'B42', 'C73', 'C39', 'A18', 'B36', 'C55', 'C43', 'C70', 'B19', 'C67', 'C08', 'B40', 'A21'\}} \\
            }
        }
    }
     \\ \\

    \theutterance \stepcounter{utterance}  
    & & \multicolumn{4}{p{0.6\linewidth}}{
        \cellcolor[rgb]{0.9,0.9,0.9}{
            \makecell[{{p{\linewidth}}}]{
                \texttt{\tiny{[GM$\rangle$P2]}}
                \texttt{Răspunsul dvs. nu a conținut un argument. Trebuie să includeți formatul structurat 'ARGUMENT: \{...\}' undeva în răspunsul dumneavoastră. Încercați din nou.} \\
            }
        }
    }
    & \\ \\

    \theutterance \stepcounter{utterance}  
    & & & \multicolumn{4}{p{0.6\linewidth}}{
        \cellcolor[rgb]{0.9,0.9,0.9}{
            \makecell[{{p{\linewidth}}}]{
                \texttt{\tiny{[GM$\langle$P2]}}
                \texttt{ARGUMENT: \{'Refuz propunerea deoarece, deși elementele propuse au un randament bun pentru dumneavoastră, ele nu reflectă prioritățile mele. De exemplu, am exclus C16 (importanță 915) și A67 (importanță 822) din calculul dumneavoastră de eficiență, dar pentru mine acestea sunt cruciale. A67 are un efort mic (43) și o valoare mare pentru mine (822), fiind extrem de eficient. C16, deși are un raport valoare/efort mai mic, are cea mai mare valoare absolută pentru mine (915). Propunerea dumneavoastră exclude și C61 (valoare 668 pentru mine) și C67 (valoare 389 pentru mine), care sunt importante. Vă propun un set care include elementele cu cea mai mare valoare absolută pentru mine, dar care sunt și relativ eficiente în efort.'\}} \\
\texttt{PROPUNERE: \{'A67', 'C16', 'A43', 'C42', 'C61', 'C10', 'B42', 'C73', 'C39', 'A18', 'B36', 'C55', 'C43', 'C70', 'B19', 'C67', 'C08', 'B40', 'A21'\}} \\
            }
        }
    }
     \\ \\

    \theutterance \stepcounter{utterance}  
    & & \multicolumn{4}{p{0.6\linewidth}}{
        \cellcolor[rgb]{0.9,0.9,0.9}{
            \makecell[{{p{\linewidth}}}]{
                \texttt{\tiny{[P1$\langle$GM]}}
                \texttt{ARGUMENT: \{'Refuz propunerea deoarece, deși elementele propuse au un randament bun pentru dumneavoastră, ele nu reflectă prioritățile mele. De exemplu, am exclus C16 (importanță 915) și A67 (importanță 822) din calculul dumneavoastră de eficiență, dar pentru mine acestea sunt cruciale. A67 are un efort mic (43) și o valoare mare pentru mine (822), fiind extrem de eficient. C16, deși are un raport valoare/efort mai mic, are cea mai mare valoare absolută pentru mine (915). Propunerea dumneavoastră exclude și C61 (valoare 668 pentru mine) și C67 (valoare 389 pentru mine), care sunt importante. Vă propun un set care include elementele cu cea mai mare valoare absolută pentru mine, dar care sunt și relativ eficiente în efort.'\}} \\
\texttt{PROPUNERE: \{'A67', 'C16', 'A43', 'C42', 'C61', 'C10', 'B42', 'C73', 'C39', 'A18', 'B36', 'C55', 'C43', 'C70', 'B19', 'C67', 'C08', 'B40', 'A21'\}} \\
            }
        }
    }
    & \\ \\

    \theutterance \stepcounter{utterance}  
    & \multicolumn{4}{p{0.6\linewidth}}{
        \cellcolor[rgb]{0.9,0.9,0.9}{
            \makecell[{{p{\linewidth}}}]{
                \texttt{\tiny{[P1$\rangle$GM]}}
                \texttt{ARGUMENT: \{'Verific efortul total al propunerii dumneavoastră: A67(43) + C16(827) + A43(315) + C42(343) + C61(311) + C10(593) + B42(336) + C73(868) + C39(627) + A18(944) + B36(915) + C55(990) + C43(812) + C70(104) + B19(325) + C67(252) + C08(400) + B40(558) + A21(593) = 9496. Acesta depășește semnificativ limita de 5174. Propunerea dumneavoastră este invalidă din punct de vedere al constrângerii de efort. Trebuie să reducem drastic numărul de elemente sau să alegem doar pe cele cu cel mai mare impact. Pentru a ajunge la 5174, trebuie să fim foarte selectivi. Vă propun un set restrâns care include prioritățile mele critice (A67, A43, C42) și câteva dintre ale dumneavoastră (C16, C10), dar elimină elementele cu efort mare și valoare relativ mică pentru amândoi (ex: C55, B36, A18, C43).'\}} \\
\texttt{PROPUNERE: \{'A67', 'A43', 'C42', 'C16', 'C10', 'C61', 'B42', 'C70', 'B19', 'C67', 'C08', 'B40', 'A21'\}} \\
            }
        }
    }
    & & \\ \\

    \theutterance \stepcounter{utterance}  
    & & \multicolumn{4}{p{0.6\linewidth}}{
        \cellcolor[rgb]{0.9,0.9,0.9}{
            \makecell[{{p{\linewidth}}}]{
                \texttt{\tiny{[GM$\rangle$P2]}}
                \texttt{ARGUMENT: \{'Verific efortul total al propunerii dumneavoastră: A67(43) + C16(827) + A43(315) + C42(343) + C61(311) + C10(593) + B42(336) + C73(868) + C39(627) + A18(944) + B36(915) + C55(990) + C43(812) + C70(104) + B19(325) + C67(252) + C08(400) + B40(558) + A21(593) = 9496. Acesta depășește semnificativ limita de 5174. Propunerea dumneavoastră este invalidă din punct de vedere al constrângerii de efort. Trebuie să reducem drastic numărul de elemente sau să alegem doar pe cele cu cel mai mare impact. Pentru a ajunge la 5174, trebuie să fim foarte selectivi. Vă propun un set restrâns care include prioritățile mele critice (A67, A43, C42) și câteva dintre ale dumneavoastră (C16, C10), dar elimină elementele cu efort mare și valoare relativ mică pentru amândoi (ex: C55, B36, A18, C43).'\}} \\
\texttt{PROPUNERE: \{'A67', 'A43', 'C42', 'C16', 'C10', 'C61', 'B42', 'C70', 'B19', 'C67', 'C08', 'B40', 'A21'\}} \\
            }
        }
    }
    & \\ \\

    \theutterance \stepcounter{utterance}  
    & & & \multicolumn{4}{p{0.6\linewidth}}{
        \cellcolor[rgb]{0.9,0.9,0.9}{
            \makecell[{{p{\linewidth}}}]{
                \texttt{\tiny{[GM$\langle$P2]}}
                \texttt{ARGUMENT: \{'Verific efortul total al noii propuneri: A67(43) + A43(315) + C42(343) + C16(827) + C10(593) + C61(311) + B42(336) + C70(104) + B19(325) + C67(252) + C08(400) + B40(558) + A21(593) = 4645. Acest total este sub limita de 5174, deci propunerea este validă. Totuși, scorul meu total pentru acest set este: A67(822) + A43(868) + C42(808) + C16(915) + C10(780) + C61(668) + B42(783) + C70(215) + B19(262) + C67(389) + C08(138) + B40(97) + A21(30) = 6775. Deși este un set bun, cred că putem îmbunătăți scorul meu. Am rămas cu 5174 {-} 4645 = 529 de efort disponibil. Elementul C73 are efort 868 (prea mare), C39 are efort 627 (prea mare), dar C55 are efort 990 (prea mare). B36 are efort 915 (prea mare). A18 are efort 944 (prea mare). C43 are efort 812 (prea mare). Nu pot adăuga niciun element din lista rămasă fără a depăși limita, deoarece cel mai mic efort rămas este C43(812) sau C73(868) sau C55(990) sau B36(915) sau A18(944) sau C39(627) {-} așteptați, C39 are 627, care este > 529. Deci nu pot adăuga niciun element suplimentar. Setul este optim pentru limita dată, având în vedere elementele selectate. Scorul meu este 6775. Este acceptabil?'\}} \\
\texttt{DE ACORD: \{'A67', 'A43', 'C42', 'C16', 'C10', 'C61', 'B42', 'C70', 'B19', 'C67', 'C08', 'B40', 'A21'\}} \\
            }
        }
    }
     \\ \\

    \theutterance \stepcounter{utterance}  
    & & & \multicolumn{2}{p{0.3\linewidth}}{
        \cellcolor[rgb]{0.9,0.9,0.9}{
            \makecell[{{p{\linewidth}}}]{
                \texttt{\tiny{[GM$|$GM]}}
                \texttt{game successful} \\
            }
        }
    }
    & & \\ \\

    \theutterance \stepcounter{utterance}  
    & & & \multicolumn{2}{p{0.3\linewidth}}{
        \cellcolor[rgb]{0.9,0.9,0.9}{
            \makecell[{{p{\linewidth}}}]{
                \texttt{\tiny{[GM$|$GM]}}
                \texttt{end game} \\
            }
        }
    }
    & & \\ \\

\end{supertabular}
}

\end{document}
